Real-world prediction tasks often include many features, most of which are only weakly relevant to the outcome. To assess our method's robustness to such features, we designed a \emph{semi-synthetic} task.
Adapting the toy generation process that produced the 2-dim. feature vectors $x$ and binary labels $y$ in Fig.~\ref{fig:toy_data_precision_90}, we append to each $x$ vector a 98-dimensional feature vector representing vitals, labs, and demographics for one randomly-chosen patient-stay in the MIMIC-III dataset \citep{johnsonMIMICIIIFreelyAccessible2016}, generating
684 train, 513 valid, and 512 test examples.
Our goal is to enforce a minimum precision of $\alpha = 0.8$. We deliberately pursue a slightly lower $\alpha$ here than in Sec.~\ref{sec:synth} to illustrate the flexibility of our methods to meet any desired precision constraint. Naturally, lowering the required minimum precision leads to better recall compared to the results in Fig.~\ref{fig:toy_data_precision_90}.
To avoid local optima, we run our method many times (25 random seeds, 2 initialization scale factors, and 2 $\lambda$ values), selecting the run that maximizes validation recall while meeting the precision constraint.

Table \ref{table:semi_synthetic_performance} shows that in this higher-dimensional problem (100 total features), only our method can meet the desired 0.8 precision even on the training set (others range from 0.63-0.73). Our method also exceeds others on test-set recall by at least 0.05. In terms of runtime, most methods are fast (15 seconds), but the AP method \citep{fathonyAPperfIncorporatingGeneric2020} shows poor scalability, requiring 2.5 hours  due to its cubic scaling. 
This poor scalability prevented our use of the AP method on the larger real-world tasks below.